\chapter{源码阅读索引}

本附录保存后续源码专题的入口。

OpenBLAS 和 BLIS 用于阅读 GEMM 的经典分层结构：block size、packing、microkernel、架构分发和多线程切分。阅读时优先找 SGEMM 路径，不要一开始陷入所有数据类型。

oneDNN 用于阅读工业深度学习后端：memory descriptor、primitive、post-op 融合、ISA 选择、卷积和矩阵乘实现。阅读时重点观察 layout 和 post-op 如何影响内核选择。

XNNPACK 用于阅读移动端和 CPU 推理内核：operator 创建、microkernel 表、indirection buffer、threadpool 和量化路径。它适合理解小算子和低延迟场景。

llama.cpp 和 ggml 用于阅读大语言模型 CPU 推理：tensor 表示、计算图、后端分发、量化 block、KV Cache、RMSNorm、RoPE、Attention 和采样。阅读时应区分模型加载、图构建、单 token decode 和 benchmark 路径。

Linux 内核源码用于理解调度、NUMA、页表、透明大页、perf 事件和内存分配。第二册不会把 Linux 内核作为主线，但性能工程遇到线程绑定、页错误、NUMA 和 perf 计数器时，需要知道这些机制来自哪里。
